% Generated draft: requires mathematical and source review.
\documentclass[11pt]{article}
\usepackage[margin=25mm]{geometry}
\usepackage{fontspec}
\setmainfont{Latin Modern Roman}
\usepackage{amsmath,amssymb,mathtools,unicode-math}
\setmathfont{Latin Modern Math}
\usepackage{enumitem,needspace,xurl,hyperref}
\setlength{\parindent}{0pt}
\setlength{\parskip}{5pt}
\setlength{\emergencystretch}{3em}
\newcommand{\N}{\mathbb{N}}
\newcommand{\Z}{\mathbb{Z}}
\newcommand{\Q}{\mathbb{Q}}
\newcommand{\R}{\mathbb{R}}
\newcommand{\C}{\mathbb{C}}
\newcommand{\F}{\mathbb{F}}
\newcommand{\A}{\mathbb{A}}
\newcommand{\PP}{\mathbb{P}}
\newcommand{\E}{\mathbb{E}}
\newcommand{\eps}{\varepsilon}
\newcommand{\dd}{\,\mathrm{d}}
\begin{document}
\section*{UnsolvedMath ID: 3800008}
\textbf{Source problem code:} catalogue entry
\subsection{Definitions and mathematical statement}
For lines $\ell_1,\ldots,\ell_n\subset\R^3$, let $d(x,\ell_i)=\inf_{y\in\ell_i}\|x-y\|_2$. Their Euclidean nearest-site Voronoi cells are
\[
V_i=\{x\in\R^3:d(x,\ell_i)\leq d(x,\ell_j)\text{ for all }j\}.
\]
The Voronoi diagram is the subdivision by nearest-line identities and ties. The third problem asks for the tight asymptotic worst-case combinatorial complexity of this diagram over configurations of $n$ lines in three-space, as $n\to\infty$.

The first two problems ask for the corresponding tight worst-case complexities of the medial surface and of an offset surface of a polyhedron with $n$ features. Under one customary solid-polyhedron convention, medial points are centers of maximal interior Euclidean balls; an outer offset at distance $r>0$ is the boundary of $P+\overline B(0,r)$. These descriptions identify common meanings, not additional source hypotheses: the record does not specify whether the medial surface is interior or exterior, whether offsets are inner or outer or offsets of the boundary, or how the offset distance is quantified.

With these conventions fixed by source review, the requested output is a matching asymptotic upper and lower bound, separately for each of the three extremal complexity functions. ``Complexity'' ordinarily counts the vertices, edges, and surface pieces of a suitable finite stratification. The precise stratification, the meaning of an input feature, and any nondegeneracy or general-position assumptions are absent here. The supplied statement reports quadratic lower-bound constructions and superquadratic upper bounds; it supplies no exact exponents or theorem hypotheses. No specific superquadratic formula is inferred.
\subsection{Short English statement}
Find tight worst-case bounds for three-dimensional line Voronoi diagrams and for medial and offset surfaces of polyhedra with $n$ features. The record reports a gap between quadratic lower bounds and superquadratic upper bounds, but leaves several input and complexity conventions unspecified.
\subsection{Sources}
\begin{itemize}
\item Local UnsolvedMath registry, permanent ID 3800008, exact source code catalogue entry, ``Medial surfaces and Voronoi diagrams of lines''; proposed by Jeff Erickson and cited proposers. Registry status partially solved, classification PARTIAL-PROGRESS.
\item Jeff Erickson, \emph{Open Problems in Computational Geometry}, source item ``Medial surfaces and Voronoi diagrams of lines'', \url{https://jeffe.cs.illinois.edu/open/}. Supplied HTML extraction accessed July 29, 2026. No individual theorem, construction, or publication citation is supplied for the reported bounds.
\item The research summary describes quadratic lower bounds and slightly superquadratic upper bounds and calls the tight bound open. The source explicitly warns that its pages have not been systematically updated since 2001; the supplied provenance therefore flags status NEEDS\_REVIEW. Public access was verified, but redistribution permission was not established and rights are also NEEDS\_REVIEW.
\end{itemize}
\end{document}
